\section*{Hoofdstuk \ref{ch:g8:colors-spectra} --- Kleuren en spectra van licht}

\begin{solution}{exo:g8:colors-spectra:1}
De aaneengesloten band kleuren die zich in wit \omterm{def:g1:light-and-shadows:source}{licht} verstopt,
uitgepakt; \omterm{prop:g8:colors-spectra:mixture}{dispersie} is het uitpakken --- elke kleur met haar
eigen afbuiging. Rood buigt het minst af, violet het meest.
\end{solution}

\begin{solution}{exo:g8:colors-spectra:2}
De regenboog vormt zich in de helft van de hemel tegenover de Zon ---
het zonlicht komt van achter de jager de regendruppels in. Elke
regendruppel speelt voor prisma (met een spiegelkaats erin).
\end{solution}

\begin{solution}{exo:g8:colors-spectra:3}
Uitpakken: het prisma waaiert wit \omterm{def:g1:light-and-shadows:source}{licht} uit tot het
\omterm{prop:g8:colors-spectra:mixture}{spectrum}. Inpakken: de draaiende regenboogschijf vervaagt tot
witachtig --- en het tweede prisma vouwt een \omterm{prop:g8:colors-spectra:mixture}{spectrum} terug tot
één witte \omterm{def:g6:light-propagation:ray}{bundel}.
\end{solution}

\begin{solution}{exo:g8:colors-spectra:4}
Gras slikt het grootste deel van het mengsel op en geeft voornamelijk
groen terug. De bladzijde geeft vrijwel alles terug --- wit; het
schoolbord geeft vrijwel niets terug en slikt de hele voorraad op ---
zwart.
\end{solution}

\begin{solution}{exo:g8:colors-spectra:5}
Wit \omterm{def:g1:light-and-shadows:source}{licht}: de sjaal geeft zijn geel terug --- gele sjaal. Blauw
\omterm{def:g1:light-and-shadows:source}{licht}: het aanbod (blauw) ontmoet het antwoord (geel) in een lege
doorsnede --- de sjaal toont zich zwart.
\end{solution}

\begin{solution}{exo:g8:colors-spectra:6}
De appel kan alleen rood teruggeven, en de blauwe schijnwerper brengt
er geen: zwart. De zakdoek geeft terug wat er aankomt --- blauw
aangeboden, geeft hij blauw terug: een blauwe zakdoek.
\end{solution}

\begin{solution}{exo:g8:colors-spectra:7}
Rood, groen, blauw. Rood $+$ groen: geel. Groen $+$ blauw: cyaan.
Blauw $+$ rood: magenta. Alle drie: wit.
\end{solution}

\begin{solution}{exo:g8:colors-spectra:8}
Alleen minuscule rode, groene en blauwe streepjes, allemaal aan ---
nergens wit. Het mengen gebeurt in je oog, dat te ver weg zit om de
streepjes te scheiden.
\end{solution}

\begin{solution}{exo:g8:colors-spectra:9}
\omterm{def:g1:light-and-shadows:source}{Lichten} tellen op wat ze meebrengen, en de volledige verzameling
kleuren samen is wit --- optellen maakt het mengsel compleet. Verven
slikken op, elk trekt haar deel van het \omterm{def:g1:light-and-shadows:source}{licht} af, en de
volledige verzameling aftrekkingen laat vrijwel niets over ---
duisternis bij commissie.
\end{solution}

\begin{solution}{exo:g8:colors-spectra:10}
De \omterm{def:g2:air-around-us:air}{lucht} steelt bij voorkeur blauw; de hele hemel ontvangt het
gestolen blauw en stuurt het omlaag door als de blauwe koepel van de
dag; de roodoranje rest overleeft de scherende weg van de
\omterm{def:g1:day-and-night:daynight}{zonsondergang}. De verduisterde \omterm{def:g2:moon-in-the-sky:moon}{Maan} wordt alleen
verlicht door \omterm{def:g1:light-and-shadows:source}{licht} dat langs de rand van de Aarde heeft
gescheerd --- rondom de hele \omterm{def:g4:solar-system:planet}{planeet} door \omterm{def:g1:day-and-night:daynight}{zonsondergangen}
gefilterd --- dus gloeit ze met dezelfde rode overlevende.
\end{solution}

\begin{solution}{exo:g8:colors-spectra:11}
De warme roodachtige lampen boden weinig groen; het echte antwoord van
het jasje --- vooral groen met wat rood dat samen een dof mengsel geeft
--- kon alleen zijn roodvriendelijke deel tonen, en dat las als
warmbruin. Het volledige aanbod van het daglicht liet het groene
antwoord spreken: bedrog door een gierig \omterm{def:g1:light-and-shadows:source}{licht}, veroordeeld door
de doorsnederegel.
\end{solution}

\begin{solution}{exo:g8:colors-spectra:12}
Letters blauw, ondergrond zwart. Onder de witte zaalverlichting geven
de blauwe letters hun blauw terug terwijl de zwarte ondergrond niets
teruggeeft: helderblauwe letters op zwart --- leesbaar. Onder de diepe
rode gloed krijgen de blauwe letters niets aangeboden dat ze kunnen
teruggeven en geeft de zwarte ondergrond zoals altijd niets terug:
gelijkmatig zwart, spandoek verdwenen. Witte letters zouden de finale
verraden: de universele teruggever antwoordt zelfs op een rode gloed
met helderrode letters, leesbaar tegen de zwarte ondergrond.
\end{solution}

\begin{solution}{pb:g8:colors-spectra:1}
\textbf{1.} Wit --- de drie primaire kleuren op volle sterkte tellen op
tot het volledige mengsel.
\textbf{2.} Geel.
\textbf{3.} Wit overhemd: geeft het aanbod terug --- geel. Blauwe
sjerp: geen blauw aangeboden --- zwart. Gele japon: het gele antwoord
past bij het gele aanbod --- schitterend geel.
\textbf{4.} Blauw achterdoek. Wit doek geeft terug wat de ontwerper
ernaartoe stuurt: één achterdoek, elke kleur op afroep --- de tegenhanger
van het witte scherm in de bioscoop.
\textbf{5.} Een groene jas: onder de rode gloed ontmoet haar groene
antwoord geen groen --- zwarte schurk; onder wit \omterm{def:g1:light-and-shadows:source}{licht} laat het
volledige aanbod het groen spreken --- rijk groen.
\textbf{6.} Een wit gewaad --- de universele teruggever: rood onder de
rode schijnwerper, groen onder groen, blauw onder blauw, wit onder alle
drie; altijd het helderste op het toneel, altijd ``witachtig'' ten
opzichte van de scène.
\textbf{7.} Elke belichting \emph{zonder} rood: de blauwe
schijnwerper, de groene, of allebei samen --- zonder rood aanbod geeft
de sjaal niets terug en \omterm{def:g2:ice-water-steam:meltfreeze}{smelt} hij samen met het zwarte gordijn.
Op het moment dat de rode schijnwerper zich erbij voegt, bevat het
aanbod precies wat de sjaal antwoordt: hij vlamt rood op, en de truc
sterft.
\textbf{8.} Goud-als-warm-geel antwoordt rood en groen gul, blauw
nauwelijks. De blauwzware stand bood weinig dat het kon teruggeven:
modderige, doffe letters. Het volledige mengsel van de middag liet het
geel oplaaien.
\textbf{9.} Rode streep: schitterend rood. Oranje en geel: gedoofde
roodachtige versies van zichzelf (hun antwoorden bevatten wat rood).
Groen: zwart. Blauw en violet: zwart (violet nog net gloeidonker). De
boog stort in tot een rood-zwarte ruïne.
\textbf{10.} De rode en de blauwe streep herleven; violet gloeit zwak
(het antwoordt blauw en wat rood); oranje en geel tonen alleen hun
zwakke rode aandeel; de groene streep blijft zwart --- magenta
\omterm{def:g1:light-and-shadows:source}{licht} bevat geen groen.
\textbf{11.} Alleen wit \omterm{def:g1:light-and-shadows:source}{licht} --- het volledige mengsel --- toont
alle zes de strepen: elke streep geeft haar eigen kleur terug, en
alleen een aanbod dat alle zes bevat, kan door alle zes worden
beantwoord. Een groene schijnwerper, hoe fel ook, verlicht één
streep en vermoordt er vijf.
\textbf{12.} Bijvoorbeeld: ``Gemengde \omterm{def:g1:light-and-shadows:source}{lichten} \emph{tellen op}:
rood, groen en blauw bouwen elke kleur en, samen, wit. En elk voorwerp
\emph{trekt af}: het geeft alleen zijn eigen aandeel terug van wat het
\omterm{def:g1:light-and-shadows:source}{licht} meebracht --- bied het niets aan dat het kan teruggeven, en
het kleedt zich in het zwart.''
\end{solution}
